\section{Related Work}
\label{sec:related}

Our work sits at the intersection of four distinct research threads: INR compression, weight-space horizontal symmetry, data-weight correspondence theory, and NTK-based analysis of INR.
We clarify the precise relationship of each to our contribution.

\subsection{INR as Data: COIN, COIN++, and Functa}

The line of work treating neural network weights as signal representations began with COIN~\citep{chen2022coin}, which showed that storing the weights of an overfitted MLP---mapping pixel coordinates to RGB values---could serve as a compressed image representation.
The key insight is that the weights \emph{are} the signal, decoupled from any architectural prior.

COIN++~\citep{dupont2022coinpp} advanced this by storing only \emph{modulations}: residual parameter shifts applied to a meta-learned base network.
This dramatically reduced encoding time and improved compression ratios.
The underlying assumption is that the meaningful variation across images lives in a low-dimensional subspace of the full parameter space.

Functa~\citep{dupont2022functa} extended this further, treating modulation vectors (and more generally, INR weights) as data points on which downstream tasks can be performed.
If image classification or retrieval can be done directly in weight space, the representation has practical value beyond compression.

Our contribution with respect to this thread is \textbf{theoretical grounding}.
We explain \emph{why} modulation vectors are low-dimensional: signals related by group transformations trace low-dimensional orbital manifolds in weight space.
A modulation vector is a coordinate along the tangent direction of this orbit.
The low-dimensionality of effective modulations is not an empirical coincidence but a geometric necessity, constrained by the dimension of the transformation group.
This theory also predicts when modulation-based approaches will degrade: when the orbit leaves the regime where the symmetry error $\delta$ is small, the low-dimensional description breaks down.

\subsection{Horizontal Symmetry: EquiGen and Permutation Equivalence}

The most directly competitive recent work is EquiGen~\citep{peignier2023equigen}, a CVPR 2025 Highlight paper.
EquiGen studies the weight-space structure of \emph{functionally equivalent} networks: different parameter vectors that represent the \emph{same} signal.
It establishes that such networks are related by a permutation group acting on hidden units, forming equivalence classes in weight space.
By projecting to an invariant latent space, EquiGen enables controlled generation of diverse but functionally equivalent networks.

EquiGen operates in the \textbf{horizontal} direction: it explores the space of weight vectors that represent a fixed signal.
It asks: given one overfitted network, what other weight vectors represent the same function?

We study the \textbf{vertical} direction: how the optimal weight vector for a signal \emph{changes} as the signal itself is transformed.
The transformation we study is not permutation of hidden units but a geometric transformation of the target data (scale, rotation).
Table~\ref{tab:horizontal_vs_vertical} makes this distinction precise.

\begin{table}[t]
    \centering
    \caption{Horizontal vs. vertical parameter-space geometry.
    The two directions study fundamentally different structures and are complementary.}
    \label{tab:horizontal_vs_vertical}
    \begin{tabular}{clll}
        \toprule
        Direction & Fixed Element & Varies Over & Representative Question \\
        \midrule
        Horizontal & Signal $f$ & Weight vector $\theta$ & What other weights represent $f$? \\
        Vertical & Weight structure & Signal family $T_g[f]$ & How does $\theta^*(f)$ move as $f$ transforms? \\
        \midrule
        \textbf{Work} & \textbf{EquiGen} & Permutation group on hidden units & Same function, equivalent networks \\
        \textbf{Ours} & Orbit geometry & Lie group on signals & Same network family, transformed signals \\
        \bottomrule
    \end{tabular}
\end{table}

The two directions are complementary: horizontal analysis helps understand the intrinsic dimensionality of the weight representation for a single function; vertical analysis characterizes how that representation changes continuously with the signal.

\subsection{Data-Weight Correspondence via Implicit Function Theorem}

A recent theoretical work (arXiv:2601.23181,~\citep{anonymous2024dataweight}) establishes a rigorous mapping between data space and weight representation space using the implicit function theorem (IFT).
The key result is that under certain regularity conditions, small changes in the data produce well-defined changes in the optimal weights, characterized by the Jacobian of the optimal-parameter mapping.

This work and ours share the IFT as a mathematical tool but apply it to different questions.
Their IFT analysis determines the \emph{conditions under which} a data-weight correspondence exists and is stable (Jacobian rank conditions).
Our IFT computation yields the specific \emph{direction} of parameter motion induced by signal transformations: the orbit tangent vector $J_X(f)$.
The two analyses are orthogonal: one concerns existence and stability of the mapping, the other characterizes its action along group directions.

Concretely, given $f \to \theta^*(f)$, arXiv:2601.23181 asks whether the mapping is smooth and invertible.
We ask: if we move $f$ along a scale or rotation transformation, what direction do $\theta^*$ move in weight space?

\subsection{NTK and Spectral Bias in INR}

The neural tangent kernel (NTK) framework~\citep{basri2020neural} characterizes the inductive bias of randomly-initialized INR.
At infinite width, the network's evolution under gradient descent is governed by the NTK, whose eigenfunctions are harmonic modes.
This explains the spectral bias of SIRENs: low-frequency components are learned faster because they correspond to larger NTK eigenvalues.

\citet{tancik2020meta} showed that MAML reshapes the empirical NTK during meta-training, steering it toward target signal structures and enabling fast adaptation.
In our framework, MAML's outer loop acts as a \emph{geometric explicitation probe}.
By maximizing gradient coherence across tasks, MAML drives the empirical NTK's leading eigenvectors to align with the true orbit tangent directions $J_X$.
This alignment accelerates the shrinkage of the symmetry error $\delta$, making the orbital structure more visible empirically.

Our \S\ref{sec:maml_alignment} provides a precise geometric interpretation of this effect, clarifying that MAML does not create orbital structure but reveals it.

\input{figures/fig_related_work_comparison}
